\documentclass[12pt,a4paper]{ctexart}
\usepackage[utf8]{inputenc}
\usepackage{amsmath, amssymb}
\usepackage{geometry}
\geometry{margin=1in}
\usepackage{enumitem}
\usepackage{array}
\usepackage{booktabs}
\usepackage{longtable}
\usepackage{multirow}
\usepackage{xcolor}
\usepackage[bf,small]{caption}
\usepackage{hyperref}
\usepackage{float}

\title{\textbf{合成高分子课堂笔记（二）}\\\large Commercialized and Engineering Polymers}
\author{Zhiheng Yi}
\date{\today}

\begin{document}

\maketitle
\tableofcontents

\section{基本概念与分类（Basic Concepts \& Classification）}

\subsection{关键定义（Key Definitions）}

\begin{itemize}[leftmargin=*]
    \item \textbf{Monomer / 单体}: A small molecule that can react chemically with other monomers to form a polymer chain. （能与其他单体反应形成聚合物链的小分子。）
    \item \textbf{Polymer / 聚合物}: A large macromolecule composed of many repeating structural units (monomers) connected by covalent bonds. （由许多重复结构单元通过共价键连接而成的大分子。）
    \item \textbf{Degree of Polymerization (DP) / 聚合度}: The number of monomer units in a polymer chain. $DP = M_n / M_0$ where $M_n$ is number-average molecular weight, $M_0$ is monomer molecular weight. （聚合物链中单体单元的数目。）
    \item \textbf{Number-average molecular weight ($M_n$) / 数均分子量}: Total weight of all polymer chains divided by the number of chains. （按聚合物链的数量平均的分子量。）
    \item \textbf{Weight-average molecular weight ($M_w$) / 重均分子量}: Weighted by the mass of chains; always larger than $M_n$ for polydisperse samples. （按链的质量加权平均的分子量，对于多分散样品通常大于$M_n$。）
    \item \textbf{Polydispersity index (PDI)}: $M_w / M_n$; indicates the breadth of molecular weight distribution. （分子量分布宽度指数。）
    \item \textbf{End group / 端基}: The functional group at the end of a polymer chain, important for characterization and further reactions. （聚合物链末端的官能团。）
    \item \textbf{Crosslinking / 交联}: Covalent bonds that connect adjacent polymer chains, forming a three-dimensional network. High crosslink density gives thermosets; low density gives elastomers. （相邻聚合物链之间的共价连接，形成三维网络。高交联密度→热固性塑料；低交联密度→弹性体。）
    \item \textbf{Thermoplastic / 热塑性塑料}: Polymers that soften when heated and harden when cooled; process is reversible; recyclable. Chains are linear or branched, not crosslinked. （加热软化、冷却硬化，过程可逆，可回收。链为线形或支化，无交联。）
    \item \textbf{Thermoset / 热固性塑料}: Polymers that undergo irreversible crosslinking during curing; once set, they cannot be remelted or reshaped. （固化过程中发生不可逆交联，一旦成型不能再熔化或重塑。）
    \item \textbf{Elastomer / 弹性体}: Polymers with low crosslinking density that can be stretched to high strains and return to original shape; rubber-like. （低交联密度聚合物，可高倍拉伸并恢复原状，类似橡胶。）
\end{itemize}

\subsection{聚合物拓扑结构（Polymer Topology）}
\begin{itemize}[leftmargin=*]
    \item \textbf{Linear polymer / 线形聚合物}: Chains without branches; chains slide easily → thermoplastic. Example: HDPE. （无支链，链易滑动→热塑性。）
    \item \textbf{Branched polymer / 支化聚合物}: Linear backbone with side chains; more heat needed to soften; still thermoplastic. Example: LDPE. （主链带有侧链，需要更多热量软化，仍为热塑性。）
    \item \textbf{Crosslinked network / 交联网络}: Adjacent chains covalently bonded → thermoset; cannot be reformed. Example: epoxy, vulcanized rubber. （相邻链共价键合→热固性，不能重塑。例：环氧树脂，硫化橡胶。）
    \item \textbf{Block copolymer / 嵌段共聚物}: A polymer consisting of long sequences (blocks) of one monomer followed by blocks of another monomer. Example: ABA triblock copolymer. （由一种单体的长链段与另一种单体的长链段交替组成。例：ABA。）
    \item \textbf{Graft copolymer / 接枝共聚物}: A backbone of one type of repeat unit with branches (grafts) consisting of another repeat unit. （一种重复单元构成主链，另一种重复单元构成侧链“接枝”。）
    \item \textbf{Random copolymer / 无规共聚物}: Two monomers randomly distributed along the chain. （两种单体沿链随机分布。）
    \item \textbf{Alternating copolymer / 交替共聚物}: Monomers alternate perfectly: ABABAB... （单体完美交替：ABABAB…）
\end{itemize}

\subsection{热性能分类（Thermal Properties）}

Thermoplastic, softens when heating, can be reshaped, recyclable. Thermoset, remains rigid, cannot reshape, non-recyclable.



\subsection{天然与合成聚合物（Natural vs Synthetic Polymers）}
\begin{itemize}[leftmargin=*]
    \item \textbf{Natural polymers / 天然聚合物}: Cotton (纤维素), starch (淀粉), natural rubber (聚异戊二烯), silk (丝蛋白), cellulose (纤维素), wool (羊毛), DNA, proteins (蛋白质).
    \item \textbf{Synthetic polymers / 合成聚合物}: Plastic bags, polyester fiber, synthetic rubber (SBR, etc.), polyurethane foams, polyethylene, nylon, PVC, PS, PET, ABS, PC.
    \item \textbf{Raw materials / 原料来源}: Crude oil or natural gas → petrochemicals (ethylene, propylene, styrene, etc.). （原油或天然气→石化产品（乙烯、丙烯、苯乙烯等）。）
\end{itemize}

\subsection{聚合物应用分类（Application-based Classification）}
\begin{itemize}[leftmargin=*]
    \item Plastic（塑料）, Fiber（纤维）, Rubber（橡胶）, Adhesive（粘合剂）, Surface coating（涂料）, Foam（泡沫）.
\end{itemize}

\section{聚合物结构层次（Levels of Polymer Structure）}

\begin{enumerate}[leftmargin=*]
    \item \textbf{Atomic level (原子级)}: Chemical structure of the repeat unit. Single bonds give flexible chains (e.g., PE); rigid aromatic rings give stiffness (e.g., Kevlar, PET). （重复单元的化学结构。单键→柔顺链；刚性芳环→刚性链。）
    \item \textbf{Molecular level (分子级)}: Topology (linear, branched, network, helical), degree of polymerization, copolymer type. （拓扑结构、聚合度、共聚物类型。）
    \item \textbf{Crystallinity (结晶度)}: 
        \begin{itemize}
            \item Amorphous (无定形): No long-range order; exhibits glass transition temperature $T_g$.
            \item Crystalline (结晶): Ordered domains; exhibits melting temperature $T_m$.
            \item Semicrystalline (半结晶): Combination of crystalline and amorphous regions. Most common.
        \end{itemize}
    \item \textbf{Phase behavior (相行为)}: How two or more polymers mix or phase separate. （多种聚合物混合时的相容性或相分离。）
\end{enumerate}

\section{大宗商品聚合物（Commodity Polymers）}

\subsection{聚烯烃（Polyolefins）}
\subsubsection{聚乙烯（Polyethylene, PE）}
Monomer: ethylene ($\mathrm{CH_2=CH_2}$). Polymerization: addition (free radical, Ziegler-Natta, metallocene).
\begin{itemize}[leftmargin=*]
    \item \textbf{LDPE (Low-Density Polyethylene / 低密度聚乙烯)}: Highly branched; low crystallinity ($\sim$40-50\%); $T_g \approx -125\,^\circ$C, $T_m \approx 105-115\,^\circ$C. Soft, flexible. Uses: squeeze bottles, film, tubing.
    \item \textbf{HDPE (High-Density Polyethylene / 高密度聚乙烯)}: Linear; high crystallinity ($\sim$70-80\%); $T_g \approx -125\,^\circ$C, $T_m \approx 125-135\,^\circ$C. Stiff, strong. Uses: milk jugs, pipes, containers.
    \item \textbf{UHMWPE (Ultra-High Molecular Weight Polyethylene / 超高分子量聚乙烯)}: Very high molecular weight (数百万 g/mol); excellent abrasion resistance; $T_m \approx 130-140\,^\circ$C. Uses: artificial joints, ballistic vests.
\end{itemize}

\subsubsection{聚丙烯（Polypropylene, PP）}
Monomer: propylene ($\mathrm{CH_3CH=CH_2}$). $T_g \approx 0\,^\circ$C (isotactic), $T_m \approx 160-170\,^\circ$C. Uses: carpet fibers, ropes, pipes, food containers, automotive parts. （等规PP的半结晶性，刚性优于PE。）

\subsubsection{聚苯乙烯（Polystyrene, PS）}
Monomer: styrene ($\mathrm{C_6H_5CH=CH_2}$). Atactic PS: amorphous, $T_g \approx 100\,^\circ$C, brittle. Uses: packaging foam (expanded PS), lighting panels, appliance components. （无规PS是无定形、脆性的。）

\subsubsection{聚氯乙烯（Polyvinyl Chloride, PVC）}
Monomer: vinyl chloride ($\mathrm{CH_2=CHCl}$). $T_g \approx 80\,^\circ$C (depends on plasticizer). Rigid PVC (unplasticized) and flexible PVC (with plasticizers). Uses: pipes, electrical wire insulation, raincoats, toys.

\subsection{树脂识别码（Resin Identification Code）}
\begin{table}[H]
\centering
\begin{tabular}{cll}
\toprule
\textbf{Code} & \textbf{Polymer} & \textbf{Recyclable / 可回收} \\
\midrule
1 & PETE (PET) & Yes \\
2 & HDPE & Yes \\
3 & PVC & Limited \\
4 & LDPE & Sometimes \\
5 & PP & Yes \\
6 & PS & Limited \\
7 & Other (PC, PLA, etc.) & Varies \\
\bottomrule
\end{tabular}
\caption{Plastic recycling codes（塑料回收代码）}
\end{table}
Note: Recyclable plastic ≠ automatically reused; proper sorting required. （可回收塑料不等于自动被再利用，需要正确分类。）


\section{工程聚合物（Engineering Polymers）}

\subsection{聚对苯二甲酸乙二醇酯（Polyethylene Terephthalate, PET）}
\begin{itemize}[leftmargin=*]
    \item Family: Polyester (ester linkage $-\mathrm{COO}-$ in main chain). （主链含酯键。）
    \item Monomers: Ethylene glycol (EG) + Terephthalic acid (TPA) or Dimethyl terephthalate (DMT). （乙二醇+对苯二甲酸或对苯二甲酸二甲酯。）
    \item Polymerization: Condensation (step-growth). （缩聚反应。）
    \item Thermal transitions: $T_g \approx 70-80\,^\circ$C, $T_m \approx 250-260\,^\circ$C. （PPT提到软化温度70$^\circ$C，即接近$T_g$。）
    \item Properties: High stiffness and strength (due to aromatic ring); semicrystalline when drawn; transparent when rapidly cooled; good gas barrier. （高刚性和强度（芳环贡献）；拉伸后半结晶；快速冷却时透明；优良的气体阻隔性。）
    \item Applications: Carbonated beverage bottles, textile fibers (polyester), food jars, Mylar films, face shields. （碳酸饮料瓶、涤纶纤维、食品罐、麦拉膜、防护面罩。）
    \item Limitation: Low $T_g$ – not for hot-filling (>70$^\circ$C). （玻璃化转变温度低，不适用于热灌装。）
\end{itemize}

\subsection{聚酰胺（Polyamide, PA / Nylon）}
\begin{itemize}[leftmargin=*]
    \item Family: Contains amide group ($-\mathrm{CONH}-$). （含酰胺键。）
    \item Natural: Proteins (wool, silk).  Synthetic: Nylon.
    \item \textbf{Nylon 6,6}:
        \begin{itemize}
            \item Monomers: Hexamethylenediamine (C6 diamine) + Adipic acid (C6 diacid). （己二胺+己二酸。）
            \item Polymerization: Condensation (eliminates water). （缩聚，失水。）
            \item Thermal transitions: $T_g \approx 50-60\,^\circ$C, $T_m \approx 260-265\,^\circ$C.
            \item Properties: High strength, toughness, abrasion resistance, good solvent resistance due to hydrogen bonding between chains. （高强度、韧性、耐磨性、良好的耐溶剂性（链间氢键）。）
            \item History: DuPont (1935); first commercial use in toothbrushes (1938); women's stockings (1940).
        \end{itemize}
    \item \textbf{Nylon 6} (from caprolactam via ring-opening polymerization): $T_g \approx 50-55\,^\circ$C, $T_m \approx 220-225\,^\circ$C.
\end{itemize}

\subsection{聚碳酸酯（Polycarbonate, PC）}
\begin{itemize}[leftmargin=*]
    \item Monomers: Bisphenol A + Phosgene (or diphenyl carbonate). （双酚A+光气。）
    \item Polymerization: Condensation (step-growth).
    \item Thermal transitions: $T_g \approx 145-150\,^\circ$C; typically amorphous (no distinct $T_m$). （无定形，仅有玻璃化转变。）
    \item Properties: Transparent, very high impact resistance (shatterproof), high refractive index (1.586) → thinner lenses. （透明、极高抗冲击性（防碎）、高折射率→镜片更薄。）
    \item Applications: Eyeglass lenses, bulletproof glass, lightweight windows, electronic components. （眼镜镜片、防弹玻璃、轻质窗户、电子元件。）
\end{itemize}

\subsection{丙烯腈-丁二烯-苯乙烯共聚物（Acrylonitrile Butadiene Styrene, ABS）}
\begin{itemize}[leftmargin=*]
    \item Composition: 15-35\% acrylonitrile, 5-30\% butadiene, 40-60\% styrene.
    \item Polymerization: Emulsion or bulk polymerization.
    \item Thermal transitions: $T_g$ of styrene-acrylonitrile (SAN) matrix $\approx 105\,^\circ$C; butadiene rubber phase provides low temperature toughness (butadiene $T_g \approx -85\,^\circ$C). Overall ABS is amorphous with $T_g \approx 105\,^\circ$C (no sharp $T_m$).
    \item Role of each component:
        \begin{itemize}
            \item Acrylonitrile (丙烯腈): Chemical resistance, heat stability.
            \item Butadiene (丁二烯): Toughness, impact strength (rubber toughening).
            \item Styrene (苯乙烯): Rigidity, processability (easy injection molding).
        \end{itemize}
    \item Properties: Opaque, impact-resistant, amorphous thermoplastic; good for 3D printing (FDM filament). （不透明、抗冲击、无定形热塑性塑料；适合3D打印。）
    \item Applications: LEGO bricks, automotive interior parts, home appliances, electronics housings. （乐高积木、汽车内饰、家电、电子设备外壳。）
\end{itemize}

\subsection{聚氨酯（Polyurethane, PU）}
\begin{itemize}[leftmargin=*]
    \item Monomers: Diisocyanate (e.g., MDI, TDI) + Polyol (diol or triol).
    \item Polymerization: Step-growth forming urethane linkage ($-\mathrm{NH-CO-O}-$).
    \item Can be thermoplastic (TPU) or thermoset (crosslinked foams).
    \item $T_g$ varies widely from $-50\,^\circ$C to $+150\,^\circ$C depending on soft/hard segment ratio.
    \item Applications: Flexible foams (mattresses), rigid foams (insulation), elastomers (wheels, seals), coatings, adhesives, Spandex fibers.
\end{itemize}

\subsection{纤维素及其衍生物（Cellulose and Derivatives）}
\begin{itemize}[leftmargin=*]
    \item Natural polymer from wood, cotton; polysaccharide of glucose; extensive hydrogen bonding → high crystallinity. （天然聚合物，来自木材、棉花；葡萄糖多糖；广泛的氢键→高结晶度。）
    \item Derivatives: Cellulose nitrate (early plastic), cellulose acetate (textiles), hydroxyethylcellulose (HEC).
    \item Hydroxyethylcellulose: Water-soluble, non-crystalline; used as thickener in shampoos to suspend dirt. （水溶性、非结晶；用作洗发水增稠剂，悬浮污垢。）
\end{itemize}


\section{高性能聚合物（High Performance Polymers）}

\begin{longtable}{p{2.8cm}p{3cm}p{3cm}p{4cm}}
\caption{High Performance Polymers / 高性能聚合物}
\tabularnewline
\toprule
\textbf{Polymer} & $\mathbf{T_g}$ (\textdegree C) & $\mathbf{T_m}$ (\textdegree C) & \textbf{Key Properties \& Applications} \\
\midrule
PEEK (聚醚醚酮) & $\sim$143 & $\sim$343 & High temp resistance ($>$250\,$^\circ$C continuous), chemical resistance; aerospace, medical implants. \\
PTFE (聚四氟乙烯) & $\sim$115 & $\sim$327 & Very low friction, non-stick (Teflon); coatings, gaskets, bearings. \\
PDMS (聚二甲基硅氧烷) & $\sim$-123 & -- (amorphous) & Silicone elastomer, biocompatible; soft lithography, contact lenses. \\
Kevlar (聚对苯二甲酰对苯二胺) & $\sim$360 (decomposes) & No $T_m$ (decomposes) & High tensile strength, aromatic rigid chains; bulletproof vests, composites. \\
PI (聚酰亚胺) & $>$250 & -- (amorphous or semicrystalline) & High thermal stability (up to 400$^\circ$C); Kapton tape, flexible circuits. \\
PVDF (聚偏氟乙烯) & $\sim$-40 & $\sim$170 & Piezoelectric, chemical resistant; sensors, membranes. \\
\bottomrule
\end{longtable}

\section{催化剂与聚合方法（Catalysts \& Polymerization Methods）}
\begin{itemize}[leftmargin=*]
    \item \textbf{Ziegler-Natta catalysts} (齐格勒-纳塔催化剂): Transition metal (Ti, Zr) + organoaluminum compound. Enable stereoregular polymerization (isotactic PP) and high molecular weight at mild conditions. Important for HDPE, isotactic PP.
    \item \textbf{Metallocene catalysts} (茂金属催化剂): Single-site catalysts giving very narrow molecular weight distribution.
    \item \textbf{Free radical polymerization} (自由基聚合): Used for LDPE (high pressure), PS, PVC.
\end{itemize}

\section{塑料回收与污染（Plastic Recycling \& Pollution）}
\begin{itemize}[leftmargin=*]
    \item Reduce single-use plastics.
    \item Recycle according to resin codes (1, 2, 5 are most recycled).
    \item For unrecyclable plastics: chemical breakdown (pyrolysis, hydrolysis, biodegradation). （对于不可回收塑料，化学分解——热解、水解、生物降解。）
    \item Bioplastics (生物塑料): From renewable resources (corn starch, cellulose). Not all bioplastics are biodegradable; some are durable (e.g., bio-PET). （来自可再生资源，并非所有生物塑料都可生物降解，有些是耐用的。）
\end{itemize}


\section*{附录1：常用聚合物中英文对照及热性能速查表}
\addcontentsline{toc}{section}{附录1：常用聚合物中英文对照及热性能速查表}
\begin{longtable}{p{2cm}p{3.5cm}p{1.8cm}p{1.8cm}p{3.5cm}}
\toprule
\textbf{Polymer} & \textbf{} & $\mathbf{T_g}$ (\textdegree C) & $\mathbf{T_m}$ (\textdegree C) & \textbf{Application Example} \\
\midrule
LDPE & 低密度聚乙烯 & -125 & 105-115 & Squeeze bottles \\
HDPE & 高密度聚乙烯 & -125 & 125-135 & Milk jugs, pipes \\
UHMWPE & 超高分子量聚乙烯 & -125 & 130-140 & Artificial joints \\
PP (isotactic) & 等规聚丙烯 & 0 & 160-170 & Carpet fibers, cups \\
PS (atactic) & 无规聚苯乙烯 & 100 & -- (amorphous) & Foam packaging \\
PVC (rigid) & 硬质聚氯乙烯 & 80 & -- (amorphous) & Pipes, window frames \\
PET & 聚对苯二甲酸乙二醇酯 & 70-80 & 250-260 & Beverage bottles, fibers \\
Nylon 6,6 & 尼龙66 & 50-60 & 260-265 & Textiles, brushes, gears \\
Nylon 6 & 尼龙6 & 50-55 & 220-225 & Carpets, fibers \\
PC & 聚碳酸酯 & 145-150 & amorphous & Eyeglass lenses, bulletproof glass \\
ABS & ABS树脂 & $\sim$105 & amorphous & LEGO bricks, electronics \\
PEEK & 聚醚醚酮 & 143 & 343 & Aerospace, medical implants \\
PTFE & 聚四氟乙烯 & 115 & 327 & Non-stick coatings \\
Kevlar & 凯夫拉（芳纶） & 360 (decomp.) & no $T_m$ & Bulletproof vests \\
PDMS & 聚二甲基硅氧烷 & -123 & amorphous & Soft contact lenses, sealants \\
PI & 聚酰亚胺 & $>$250 & varies & Kapton tape, flexible PCBs \\
\bottomrule
\caption{Comprehensive Polymer Datasheet}
\end{longtable}

\section*{附录2：中英文术语对照表}
\addcontentsline{toc}{section}{附录2：中英文术语对照表}
\begin{longtable}{p{10cm}p{5cm}}
\toprule
\textbf{English} & \textbf{中文} \\
\midrule
Addition polymerization & 加聚反应 \\
Condensation polymerization & 缩聚反应 \\
Degree of polymerization (DP) & 聚合度 \\
Number-average molecular weight ($M_n$) & 数均分子量 \\
Weight-average molecular weight ($M_w$) & 重均分子量 \\
Polydispersity index (PDI) & 多分散指数 \\
End group & 端基 \\
Crosslinking & 交联 \\
Thermoplastic & 热塑性塑料 \\
Thermoset & 热固性塑料 \\
Elastomer & 弹性体 \\
Linear polymer & 线形聚合物 \\
Branched polymer & 支化聚合物 \\
Block copolymer & 嵌段共聚物 \\
Graft copolymer & 接枝共聚物 \\
Random copolymer & 无规共聚物 \\
Alternating copolymer & 交替共聚物 \\
Amorphous & 无定形的 \\
Crystalline & 结晶的 \\
Semicrystalline & 半结晶的 \\
Glass transition temperature ($T_g$) & 玻璃化转变温度 \\
Melting temperature ($T_m$) & 熔融温度 \\
Polyethylene (PE) & 聚乙烯 \\
Low-density PE (LDPE) & 低密度聚乙烯 \\
High-density PE (HDPE) & 高密度聚乙烯 \\
Ultra-high molecular weight PE (UHMWPE) & 超高分子量聚乙烯 \\
Polypropylene (PP) & 聚丙烯 \\
Polystyrene (PS) & 聚苯乙烯 \\
Polyvinyl chloride (PVC) & 聚氯乙烯 \\
Polyethylene terephthalate (PET) & 聚对苯二甲酸乙二醇酯 \\
Polyamide (PA) / Nylon & 聚酰胺 / 尼龙 \\
Nylon 6,6 & 尼龙66 \\
Polycarbonate (PC) & 聚碳酸酯 \\
Acrylonitrile butadiene styrene (ABS) & 丙烯腈-丁二烯-苯乙烯共聚物 \\
Polyurethane (PU) & 聚氨酯 \\
Polyether ether ketone (PEEK) & 聚醚醚酮 \\
Polytetrafluoroethylene (PTFE) & 聚四氟乙烯 \\
Polydimethylsiloxane (PDMS) & 聚二甲基硅氧烷 \\
Kevlar & 凯夫拉（芳纶） \\
Polyimide (PI) & 聚酰亚胺 \\
Polyvinylidene fluoride (PVDF) & 聚偏氟乙烯 \\
Resin identification code & 树脂识别码 \\
Bioplastic & 生物塑料 \\
Cellulose & 纤维素 \\
Ziegler-Natta catalyst & 齐格勒-纳塔催化剂 \\
Metallocene catalyst & 茂金属催化剂 \\
Free radical polymerization & 自由基聚合 \\
Emulsion polymerization & 乳液聚合 \\
Interfacial polymerization & 界面聚合 \\
\bottomrule
\end{longtable}

\newpage
\section*{重点提示}
\begin{enumerate}[leftmargin=*]
    \item \textbf{区分热塑性与热固性}: 本质是线形/支化（可逆） vs 交联网络（不可逆）。
    \item \textbf{写出聚合物的重复单元结构}: 例如 PET 的重复单元为 $-\mathrm{CO}-\mathrm{C_6H_4}-\mathrm{COO}-\mathrm{CH_2CH_2O}-$。
    \item \textbf{解释结构与性能关系}: 芳环 → 刚性/高 $T_m$；氢键（尼龙）→ 强度；交联 → 弹性或热固性。
    \item \textbf{树脂代码 1-7}: 1(PET), 2(HDPE), 3(PVC), 4(LDPE), 5(PP), 6(PS), 7(其他)。
    \item \textbf{热转变数据}:  PET ($T_g\sim70^\circ$C) 不适合热灌装；PC 的 $T_g$ 高，可用作耐热透明材料；PP 的 $T_m$ 高于 PE。
    \item \textbf{尼龙合成}: 尼龙66为缩聚，尼龙6为开环聚合。
    \item \textbf{ABS三组分功能}: 丙烯腈 (耐化学品/热稳定性)、丁二烯 (韧性)、苯乙烯 (刚性/加工性)。
\end{enumerate}

\end{document}
